\documentclass{article}
\usepackage{iclr2027_conference,times}

% math_commands.tex from the ICLR template is not loaded: the paper uses none of
% its 427 macros, and \eqref comes from amsmath below.

\usepackage{hyperref}
\usepackage{url}
\usepackage{graphicx}
\usepackage{amsmath,amssymb}
\usepackage{booktabs}
\usepackage{subcaption}
\usepackage{xcolor}

% ---------------------------------------------------------------------------
% Notation used throughout: the two sectors of an agent's state, the fields
% that drive learning, and the order parameters.
% ---------------------------------------------------------------------------
\newcommand{\hw}{h_{w}}
\newcommand{\hmu}{h_{\mu}}
\newcommand{\Fw}{F_{w}}
\newcommand{\Fmu}{F_{\mu}}
\newcommand{\FC}{F_{C}}
\newcommand{\FV}{F_{V}}
\newcommand{\Rwmu}{R_{w,\mu}}
\newcommand{\Rmuc}{R_{\mu,c}}
\newcommand{\Rcw}{R_{c,w}}
\newcommand{\BI}{B_{I}}
\newcommand{\BA}{B_{A}}
\newcommand{\fd}{f_{d}}

% The ICLR style sets titles in \LARGE; at that size this title wraps onto a
% second line and hyphenates ("Neural Net-works"). One step down fits it on one
% line. \texorpdfstring keeps the size command out of the PDF metadata string.
\title{Polarization and Discrimination \\ in a Multi-Agent System of Neural Networks}

\author{Anonymous authors\\Paper under double-blind review}

%\iclrfinalcopy % Uncomment for camera-ready, NOT for submission.

\begin{document}

\maketitle

\begin{abstract}
Populations of learning agents are increasingly deployed as societies: agents that
read each other's outputs, weigh how far to trust them, and revise their positions
accordingly. We ask whether discriminatory group structure can emerge in such a
population from the \emph{learning rule alone}, with no biased data and no
group-level preference anywhere in the system. We study a society of adaptive
agents that each carry a classifier and update, online and optimally, both their
opinions and the trust they assign to every other agent, and inject the minimal
possible bias: a fraction $\fd$ of agents extend their representation of a message
with a feature that carries no information about its content and depends only on
the sender's group label, at strength $d$. The result is a phase diagram in
$(d,\fd)$ with four regimes separated by sharp transitions --- neutral polarisation
unrelated to group membership, two discriminatory phases distinguished by whether
hostility still needs disagreement to sustain it, and, for bias favouring the
out-group, a frustrated spin-glass regime in which no coherent faction can form.
Three findings bear on how such systems should be built and evaluated.
Discrimination requires a quorum of biased agents --- below $\fd\approx0.4$ the
population is neutral however strong the bias --- and past it the transition is
abrupt, so a population can be discriminatory while every agent in it passes an
individual audit. The effect acts entirely through the sector of the algorithm
that estimates how reliable other agents are, so it is a hazard specific to agents
that model each other rather than a generic consequence of an irrelevant feature.
And the \emph{order} of polarisation is controlled by the complexity of the agenda:
below $\alpha = P/K \approx 1.7$ distrust forms first and drags opinions after it,
above it opinions separate first. The order parameters we use are defined only
through pairwise opinion alignment and pairwise trust, so they are measurable in
any population of interacting agents, whatever the agents are made of.
\end{abstract}

\section{Introduction}
\label{sec:intro}

Systems built from many interacting learned agents are now routine: agents that
debate, review each other's work, negotiate, and simulate populations
\citep{park2023generative,du2023debate,chuang2024simulating}. Each agent reads
what others produce, forms some judgement of how much to rely on them, and adjusts
its own position accordingly. That is a society, and societies have collective
states that no individual chose. The question this paper asks is whether one of the
collective states available to such a system is \emph{discrimination}: a
population that sorts itself into mutually distrustful groups along a label that
carries no information.

The usual explanation for group bias in machine learning locates it in the data:
a model trained on a biased corpus reproduces the bias. That account cannot
explain a population-level effect, because it says nothing about interaction. We
therefore ask a different question. Suppose every agent is individually
well-specified --- learning by an algorithm that is provably optimal in the
setting it was designed for --- and suppose the only defect is that some agents
attend to one irrelevant feature of a message, namely who sent it. What does the
\emph{population} do?

We study this in a model in which the analysis can be carried out: a society of
$N$ agents, each carrying a perceptron over a $K$-dimensional embedding of
issues, interacting by stating opinions about issues drawn from a shared agenda
of size $P$. Two features make the model unusually apt for the question. First,
agents can agree on some issues and disagree on others, so opinion distance
is graded rather than binary. Second, learning is derived rather than posited:
the update follows from a Bayesian step on the agent's belief, projected back
onto a Gaussian family by maximum entropy
\citep{opper1996online,opper1998bayesian,caticha2020entropic}, and it carries
alongside the weights an explicit, dynamical estimate of how noisy each other
agent is as a channel --- a quantity we can read as trust
\citep{alves2016distrust,caticha2011agent}. Cognitive-sounding behaviour then
appears without being written in: surprise, blame attribution between
belief-revision and source-derision, and the ability to learn the \emph{opposite}
of what a distrusted source says.

Into this society we inject the smallest bias that can be correlated with a group
label. A fraction $\fd$ of agents shift their perceived agreement with a message
by $\pm d$ according to whether the sender shares their label. Nothing else
distinguishes the groups; the labels are informationally inert; the two groups are
treated symmetrically, so any asymmetry that emerges is spontaneous. The
resulting collective behaviour, as a function of $(d,\fd)$, is the subject of the
paper.

\paragraph{Contributions.}
\begin{itemize}\itemsep1pt
  \item \textbf{A phase diagram of emergent discrimination}
    (\S\ref{sec:results}): four regimes with sharp boundaries --- neutral
    polarisation unrelated to class, two discriminatory phases distinguished by
    whether opinion still tracks the group split, and a frustrated spin-glass
    regime when the bias favours the out-group. Discrimination needs a quorum:
    below $\fd\approx0.4$ the society stays neutral however strong the bias, and
    above it the transition in $d$ sharpens rapidly.
  \item \textbf{A mechanism specific to good learners}
    (\S\ref{sec:discrimination}): the bias displaces the separatrix of the
    learning flow, enlarging one dissonance-free basin at the other's expense.
    It acts only through the algorithm's trust sector, so optimising for the
    small-group setting is what makes a population vulnerable in the large-group
    one.
  \item \textbf{Which comes first, distrust or disagreement}
    (\S\ref{sec:no-disc}): even with no bias at all, the order of polarisation is
    set by the agenda
    complexity $\alpha=P/K$, reversing at $\alpha\approx1.7$. This turns a
    contested question about polarisation
    \citep{fiorina2008polarization,iyengar2012affect} into a statement about a
    measurable parameter.
  \item \textbf{Order parameters that do not depend on the substrate}
    (\S\ref{sec:order-params}). They need only, per pair of agents, a signed
    opinion alignment and a signed trust --- never weights, gradients or the
    embedding dimension. Two elicited matrices per population therefore suffice to
    locate any interacting multi-agent system on the diagram, which matters because
    the per-agent audits that are standard practice cannot see the phase at all.
\end{itemize}

Our reading is deliberately narrow: these are simple machines, and extending
conclusions from them to human societies is not automatic. What the model does
support is a claim about learning systems --- that an inference rule optimal for a
pair of agents, transplanted into a large population, can convert an
informationally worthless label into a stable architecture of group distrust. That
is a failure mode of the interaction rather than of the training data, and it is
invisible to any evaluation that examines agents one at a time.
\section{Related work}
\label{sec:related}

\paragraph{Societies of language-model agents.} Multi-agent language systems have
been used to simulate populations and to study what emerges from their
interaction: generative agents with memory and social behaviour
\citep{park2022social,park2023generative}, debate procedures that improve
factuality \citep{du2023debate}, opinion exchange over networks
\citep{chuang2024simulating}, and language models as stand-ins for survey
respondents \citep{argyle2023out,aher2023using}, with systematic biases documented
in simulated debate \citep{taubenfeld2024systematic}. This literature is largely
empirical: it reports what a particular population does. What it lacks is a model
in which the collective outcome can be \emph{derived} and its phase boundaries
located --- which is what we supply, along with order parameters cheap enough to
measure on any such population.

\paragraph{Bias and fairness in machine learning.} The dominant framing locates
group bias in data and in single-model decisions, as representational harm
\citep{blodgett2020language,bender2021dangers} or as allocative unfairness with
formal criteria attached
\citep{dwork2012fairness,hardt2016equality,barocas2023fairness}. Our mechanism is
orthogonal and complementary: the outcome here is produced by an irrelevant
feature interacting with the \emph{learning rule} across a group of agents, and it
is undetectable by any audit that inspects agents one at a time, since a single
agent's bias is $O(d)$ while the collective state it induces is discontinuous in
$d$. That irrelevant features degrade a classifier is old; that they can
restructure a \emph{society} of classifiers is the point here.

\paragraph{Opinion dynamics and social balance.} Models of consensus and
polarisation in interacting populations run from bounded-confidence and mixing
rules \citep{deffuant2000mixing,hegselmann2002opinion} through culture
dissemination \citep{axelrod1997dissemination} to reinforcement-driven echo
chambers \citep{baumann2020modeling}; see \citet{castellano2009statistical} for a
survey. In most of that work an agent's state is a scalar or a discrete label, so
two agents either agree or do not. Representing agents as classifiers over an
issue space \citep{caticha2011agent} makes partial agreement possible, which is
what lets us ask whether distrust follows disagreement or the reverse. Our second
diagnostic, social balance over triples
\citep{heider1946attitudes,cartwright1956structural,antal2005dynamics,marvel2011continuous},
is usually applied to sign structures that are themselves the primitive objects;
here both sign structures are generated by inference, and their balance is what
separates organised bloc formation from disorder --- the distinction that
identifies the frustrated phase, which pair correlations alone would report as
nothing happening.

\paragraph{Optimal on-line learning, and affective polarisation.} The learning rule
is the optimised on-line algorithm for the student--teacher problem, given a Bayesian
reading by \citet{opper1996online} and an entropic-dynamics formulation by
\citet{caticha2020entropic}; the identification of the inferred channel-noise level
with distrust, and its promotion to a dynamical variable, is due to
\citet{alves2016distrust}. Appendix~\ref{app:genealogy} gives the fuller genealogy.
Our contribution is not the algorithm but what it does to a population when one
class-correlated feature is added. Separately, whether groups come to dislike each
other because they disagree or the reverse is contested in political science
\citep{fiorina2008polarization,iyengar2012affect,iyengar2015fear}; we do not
adjudicate it, but identify a parameter that selects which occurs.


\section{A multi-agent system of neural networks that learn from each other}
\label{sec:model}

\subsection{Microscopic variables}
\label{sec:micro}

A society of $N$ agents discusses an agenda of $P$ issues. An issue is a vector
$\hat{x}\in\mathbb{R}^{K}$: the embedding of an assertion in the $K$ dimensions
that are relevant to the question at hand. All agents embed issues identically and
differ only in what they make of them, so the model separates the representation of
an issue from the opinion an agent holds about it.

Each agent carries a single-layer perceptron with weights $w_I\in\mathbb{R}^{K}$
and states the opinion $\sigma_I = \mathrm{sign}(w_I\cdot\hat{x})\in\{\pm 1\}$.
This is the least structure that lets two agents agree on some issues and disagree
on others; a model in which each agent is a single spin cannot do that, and the
possibility of partial agreement turns out to drive the collective behaviour we
report. Agents also belong to one of two classes, $A$ or $B$, recorded as
$\kappa_I=\pm1$: a label that is visible to others and carries no information about
any issue. Until \S\ref{sec:discrimination} nothing in the dynamics refers to it.

An agent's state has two sectors. The \emph{opinion} sector is a Gaussian belief
over its own weights, with mean $\hat{w}_I$ and covariance $C_I$. The
\emph{trust-attribution} sector is a Gaussian belief, for every other agent $J$,
over how unreliable $J$ is as a source: mean $\mu_{J|I}$ (the \emph{distrust} $I$
assigns to $J$) and variance $V_{J|I}$. We call it the trust sector for short. Writing the joint belief as a product of Gaussians,
%
\begin{equation}
  Q(w_I, z_I \mid \lambda_I) \;=\; \mathcal{N}(\hat{w}_I, C_I)
  \prod_{J\neq I}\mathcal{N}(\mu_{J|I}, V_{J|I}),
  \label{eq:belief}
\end{equation}
%
the microscopic variables of the society are
$\lambda_I=\{\hat{w}_I, C_I, \mu_{J|I}, V_{J|I}\}$. The means fix what an agent
thinks and whom it trusts; the (co)variances fix how firmly.

\paragraph{Certainty and immobility are the same thing.} The second moments play a double role worth naming before the update is written down. $C_I$ and $V_{J|I}$ are the agent's uncertainties, but they also multiply every change it makes. They are effective, adaptive learning rates as much as error bars, and in the limit $C_I, V_{J|I}\to0$ the dynamics freezes: an agent absolutely certain of its opinions and of its trust in everyone cannot be moved by any message, however surprising. That is the Bayesian statement that evidence cannot shift a delta-function prior.

\subsection{Dynamics}
\label{sec:dynamics}

At each step an issue $\hat{x}$, an emitter $e$ and a receiver $r$ are drawn
uniformly. The emitter states $\sigma_e$; the receiver treats it as a noisy report
of an unknown true label, where the noise level is precisely what its distrust of
the emitter encodes. A Bayesian step on \eqref{eq:belief} leaves the exponential
family, and projecting the posterior back onto a Gaussian by maximum entropy ---
relative to the prior, subject to the first two moments --- keeps the belief in the
form \eqref{eq:belief} and yields a closed update. This is the entropic-dynamics
construction, and it is what makes the algorithm optimal in the small-group
(student--teacher) setting rather than merely plausible.
Appendix~\ref{app:derivation} carries the derivation, including the two modelling
choices that keep the algebra Gaussian: taking the probability that the emitter
conveys the wrong label to be $\Phi(z)$ with $z$ Gaussian, and letting an
inexperienced agent be correspondingly unsure it has parsed the issue as the
emitter did.

Two scaled fields carry everything the update depends on:
%
\begin{equation}
  \hw \;=\; \frac{\hat{w}_r\cdot\hat{x}\,\sigma_e}{\gamma_C},
  \qquad
  \hmu \;=\; \frac{\mu_{e|r}}{\gamma_V},
  \qquad
  \gamma_C=\sqrt{1+\hat{x}\cdot C_r\hat{x}},\quad \gamma_V=\sqrt{1+V_{e|r}} .
  \label{eq:fields}
\end{equation}
%
We call $\hw$ the \emph{opinion field}: it is positive when the receiver already
agrees with the message and large when it holds that opinion confidently. We call
$\hmu$ the \emph{distrust field}: positive when the receiver distrusts the emitter.
In these variables the evidence for the message is
%
\begin{equation}
  Z \;=\; \Phi(\hw) + \Phi(\hmu) - 2\,\Phi(\hw)\Phi(\hmu),
  \label{eq:evidence}
\end{equation}
%
which is small in two quadrants --- agreeing with a distrusted emitter, and
disagreeing with a trusted one --- and large in the other two. Low evidence is
surprise, and surprise is what drives learning. The update is a gradient step in
$\log Z$ in both sectors,
%
\begin{align}
  \hat{w}_r &\leftarrow \hat{w}_r + \frac{\Fw}{\gamma_C}\,\sigma_e\,C_r\hat{x},
  &
  C_r &\leftarrow C_r + \frac{\FC}{\gamma_C^{2}}\,(C_r\hat{x})(C_r\hat{x})^{\!\top},
  \label{eq:update-ideo}\\
  \mu_{e|r} &\leftarrow \mu_{e|r} + \frac{\Fmu}{\gamma_V}\,V_{e|r},
  &
  V_{e|r} &\leftarrow V_{e|r} + \frac{\FV}{\gamma_V^{2}}\,V_{e|r}^{2},
  \label{eq:update-aff}
\end{align}
%
with the four modulation functions given by the derivatives of $\log Z$:
%
\begin{align}
  \Fw &= \bigl(1-2\Phi(\hmu)\bigr)\frac{g(\hw)}{Z},
  &\Fmu &= \bigl(1-2\Phi(\hw)\bigr)\frac{g(\hmu)}{Z},
  \nonumber\\
  \FC &= -\Fw(\Fw+\hw),
  &\FV &= -\Fmu(\Fmu+\hmu),
  \label{eq:modulation}
\end{align}
%
where $g$ is the standard normal density. The first update in
\eqref{eq:update-ideo} is a Hebbian step $\sigma_e\hat{x}$ with a tensorial,
self-annealing learning rate --- the $C_r$ of \S\ref{sec:micro} in its
learning-rate role.

Three properties of \eqref{eq:modulation} do the work in what follows, and none of
them was put in by hand. \textbf{Trust gates learning and can invert it}: the
prefactor $1-2\Phi(\hmu)$ in $\Fw$ is negative when the emitter is distrusted, so
a distrusted emitter is not ignored but \emph{anti-learned} from, generating
effectively antiferromagnetic couplings. (A comparable reversal is documented in
humans, where a norm advertised by a distrusted source can push behaviour the
other way, \citealp{keizer2011reversal}.) \textbf{Agreement builds trust}:
symmetrically, $1-2\Phi(\hw)$ in $\Fmu$ is negative when the receiver agrees, so
whether $\mu$ rises or falls is decided by the \emph{sign of the perceived
agreement} --- the hinge that \S\ref{sec:discrimination} turns.
\textbf{Only one sector yields at a time}: the sectors are mirror images,
$\Fw(x,y)=\Fmu(y,x)$ and $\FC(x,y)=\FV(y,x)$, and the surprise is absorbed by
whichever field is smaller in magnitude, crossing over at $\hmu=\hw$
(Figure~\ref{fig:modulation}). Meeting a distrusted agent who agrees with you costs
you either your distrust or your opinion, whichever you hold less firmly --- blame
attribution emerging from the inference problem rather than assumed.

\begin{figure}[t]
  \centering
  \includegraphics[width=0.86\linewidth]{figures/modulation_contours}
  \caption{\textbf{Learning is driven by dissonance.} The two modulation functions
  that move the means, over the plane of the opinion field $\hw$ and the distrust
  field $\hmu$: $\Fw$, which revises the receiver's opinion, and $\Fmu$, which
  revises its trust in the emitter. Nothing happens where the message is
  unsurprising; everything happens in the two dissonant quadrants, agreeing with a
  distrusted emitter (upper right) and disagreeing with a trusted one (lower left).
  The two panels are reflections of each other across the green line $\hmu=\hw$,
  which is the symmetry $\Fw(x,y)=\Fmu(y,x)$ between the sectors, and that line is
  also where blame for a surprise passes from one sector to the other
  (Appendix~\ref{app:slices}). The two functions that anneal the uncertainties,
  $\FC$ and $\FV$, obey the same symmetry and are shown in
  Appendix~\ref{app:modulation}. Both functions diverge where the evidence goes to
  zero in the far corners, so the colour scale is clipped at $\pm3.2$.}
  \label{fig:modulation}
\end{figure}

\section{Emergent Polarisation}
\label{sec:no-disc}

Nothing in \S\ref{sec:model} refers to an agent's class, so before adding any bias
it is worth asking what this society does unaided. It polarises, spontaneously and
completely --- and the \emph{order} in which it polarises turns out to be controlled
by a single ratio.

\subsection{The society splits in two}
\label{sec:splits}

Figure~\ref{fig:polarisation} shows what happens to a society of $N=40$ agents left
to itself at $d=0$. Nothing in the figure involves a class label or any order
parameter; it is the raw state of the society at the start and at the end.

Two questions have to be separated. Has every pair of agents taken a side? And do
those sides amount to exactly two camps, rather than several? The distributions in
Figure~\ref{fig:polarisation} answer the first; a count over triples answers the
second. The trust
matrix is the starker of the two: it saturates, every in-faction pair at
$\eta\approx+1$ and every cross-faction pair at $\eta\approx-1$.

Both panels are distributions over pairs, so they do not depend on how the agents
are labelled and cannot be an artefact of any ordering imposed on them. Opinion
overlap begins unimodal and centred on zero --- width $0.178$
against the $1/\sqrt{K}=0.183$ expected of random vectors in $K=30$ dimensions, with
$40\%$ of pairs inside $|\rho|<0.1$ --- and ends bimodal with peaks at $-0.59$ and
$+0.63$ and \emph{no} pair left inside $|\rho|<0.1$. Trust begins spread across the
interior, nowhere near saturation, and ends with $35\%$ of ordered pairs beyond
$|\eta|>0.99$ and a mean $|\eta|$ of $0.986$. Agents have neither converged on a
consensus nor stayed independent, and the saturation of the trust matrix is something
the dynamics builds rather than something the initial condition supplied.

That the sides are two rather than many needs a different measurement, because a
bimodal distribution is equally consistent with three or four mutually hostile
groups. The structure theorem for signed complete graphs
\citep{cartwright1956structural} supplies one: such a graph is balanced --- every
triple has a positive product of signs --- exactly when its vertices split into two
mutually hostile cliques. Counting the fraction of balanced triples, which discards
the magnitudes that $\BI$ and $\BA$ weight by, therefore counts factions.
Three equal groups would give $3/4$ and random signs $1/2$. The eight societies of
Figure~\ref{fig:polarisation} contain $84\,694\,400$ triples between them. In trust,
\emph{not one} is unbalanced. In opinion, six of the eight societies are likewise
perfectly balanced and the other two hold a residue of $1\,590$ triples in total ---
two thousandths of one per cent --- which is traceable to the $29$ pairs, of
$638\,400$, whose overlap is smaller than $0.05$ in magnitude and whose sign is
therefore close to arbitrary. The society has two factions, not three or more.

They are also the \emph{same} two factions in both sectors: agents end up trusting
those who agree with them and distrusting those who do not, which is why the two
panels of Figure~\ref{fig:polarisation} split at the same moment. Nothing in
\S\ref{sec:model} imposes that alignment; it is what the sign structure of the
modulation functions produces, and \S\ref{sec:order-params} puts a number on it.

Two things about that split matter for what follows. It is \emph{spontaneous}: which
agent lands in which faction is decided by the initial condition, not by anything
intrinsic, and the two are equal only up to fluctuation --- the eight societies of
Figure~\ref{fig:polarisation} split between $190/210$ and $212/188$ of their $400$
agents, against $24$ against $16$ in an otherwise identical society of $N=40$, the
relative imbalance falling roughly as $N^{-1/2}$. And it is
\emph{class-blind}: the labels $\kappa_I$ exist, but the dynamics has not touched
them, so the factions cut across them. Everything \S\ref{sec:discrimination} adds is
a mechanism for aligning this pre-existing split with the labels.

\begin{figure}[t]
  \centering
  \includegraphics[width=1.0\linewidth]{figures/polarisation}
\caption{\textbf{An unbiased society splits in two.} Eight independent societies
  of $N=400$ agents at $d=0$, simple agenda, before and after the calibrated $500$
  interactions per ordered pair, with their pairs pooled: $638\,400$ unordered pairs
  for the opinion overlap and $1\,276\,800$ ordered pairs for the trust. Blue is the
  initial state, red the final one. Neither panel depends on how the agents are
  labelled, so neither can be an artefact of sorting them. Opinion overlap starts unimodal at zero, of width $1/\sqrt{K}$ as
  random vectors in $K=30$ dimensions require, and ends bimodal; trust starts spread
  across the interior and ends saturated at $\pm1$, so the split in trust is built by
  the dynamics and not supplied by the initial condition. Bimodality alone would show
  only that every pair has taken a side; that the sides are exactly \emph{two} blocs
  is the sign-balance figure quoted in the text.}
  \label{fig:polarisation}
\end{figure}

\subsection{Order parameters}
\label{sec:order-params}

A society's state is a set of weight vectors and an $N\times N$ matrix of distrust.
To compare states across parameters --- and, in principle, across populations built
from different kinds of agent --- we need statistics that depend on neither the
dimension $K$ nor the parameterisation.

Two pairwise quantities do the work. The \emph{opinion overlap} of two agents is
the cosine of their weight vectors,
$\rho_{IJ}=\hat{w}_I\cdot\hat{w}_J/(\|\hat{w}_I\|\|\hat{w}_J\|)$, which is $+1$
when they agree on every issue and $-1$ when they disagree on every one. The
\emph{trust} a receiver $I$ places in an emitter $J$ is $\eta_{J|I}=1-2\Phi(\hmu)$,
which is $+1$ when $J$ is fully trusted and $-1$ when fully distrusted: the distrust
field mapped onto a bounded, sign-carrying scale. Their correlation,
%
\begin{equation}
  \Rwmu = \bigl\langle (\eta_{I|J}+\eta_{J|I})\,\rho_{IJ} \bigr\rangle,
  \label{eq:rwmu}
\end{equation}
%
averaged over unordered pairs and normalised to $[-1,1]$, asks whether agents trust
those who agree with them. It is the single number behind the alignment of the two
panels of Figure~\ref{fig:polarisation}: zero at initialisation, where opinion and
trust are independent by construction, and $+0.56$ in the polarised society at
$\alpha=0.17$. How far it can rise is capped by the agenda, and
\S\ref{sec:no-disc} returns to why: at $\alpha=3.3$ the same measurement gives
$+0.93$.

A correlation cannot distinguish a society split into two coherent camps from one
in which distrust is spread inconsistently, so we also measure social balance in the
sense of \citet{heider1946attitudes} and \citet{cartwright1956structural}. A triple
of agents is balanced in opinion when $b^{I}_{IJK}=\rho_{IJ}\rho_{JK}\rho_{KI}>0$
and balanced in trust when
$b^{A}_{IJK}=\tfrac{1}{2}(\eta_{IJ}\eta_{JK}\eta_{KI}+\eta_{JI}\eta_{IK}\eta_{KJ})>0$;
the aggregates
%
\begin{equation}
  \BI = \bigl\langle b^{I}_{IJK}\bigr\rangle_{(IJK)},
  \qquad
  \BA = \bigl\langle b^{A}_{IJK}\bigr\rangle_{(IJK)}
  \label{eq:balance}
\end{equation}
%
run over all $\binom{N}{3}$ triples. (The subscripts are historical, from the older
names for the two sectors --- ideological and affective --- and do not index an
agent.) A randomly
initialised society has $\BI\approx\BA\approx0$: half its triples are frustrated. A
society that has settled into two internally coherent, mutually opposed blocs has
$\BI=\BA=1$, whatever the blocs are made of --- and a society that cannot settle at
all stays near zero or goes negative. Balance therefore separates \emph{organised}
disagreement from mere disorder, which is the distinction the phase diagram of
\S\ref{sec:results} turns on. Both aggregates are computed in closed form from
matrix products rather than by enumerating triples
(Appendix~\ref{app:closed-form}).

Note that none of \eqref{eq:rwmu} or \eqref{eq:balance} mentions weights, gradients
or $K$. They need only, per pair of agents, a signed alignment of opinion and a
signed trust. Any population in which those two things can be elicited --- whatever
its agents are internally --- can therefore be placed on the same axes, and two
$N\times N$ matrices is a cheap enough measurement to run on a deployed system.
\S\ref{sec:discussion} returns to why a per-agent audit cannot substitute for it.

\subsection{Which polarisation comes first}

The split of \S\ref{sec:splits} is what the modulation functions make inevitable:
learning drives each agent away from the dissonant quadrants of
Figure~\ref{fig:modulation}, and the two consonant quadrants --- agree-and-trust,
disagree-and-distrust --- are equally attractive. Both balances therefore rise from
zero towards one; frustration is annealed away and what remains is structured
opposition. Which balance rises first is the interesting part. Figure~\ref{fig:trajectories}
follows the society through the $(\BI,\BA)$ plane from its random initial condition,
at $d=0$, for agendas spanning four orders of magnitude in the ratio
%
\begin{equation}
  \alpha \;=\; \frac{P}{K},
  \label{eq:alpha}
\end{equation}
%
the number of issues in play per dimension of the space they are represented in.
For a simple agenda the path runs up the $\BA$ axis: trust structure forms
almost completely before opinions begin to sort, ending at
$(\BI,\BA)=(0.02,0.97)$ for $\alpha=0.03$. For a complex agenda the path stays
below the diagonal, ending at $(0.97,0.74)$ for $\alpha=333$: opinions separate
first and distrust follows. The crossover is at $\alpha\approx1.7$, where the
trajectory lies along the diagonal.

\begin{figure}[t]
  \centering
  \includegraphics[width=0.60\linewidth]{figures/agenda_trajectories}
  \caption{\textbf{Which polarisation comes first is set by the agenda.}
  Trajectories in the $(\BI,\BA)$ plane with no discrimination ($d=0$), one curve
  per agenda complexity $\alpha=P/K$, markers equally spaced in time --- they crowd
  at the end because the annealed dynamics slows, as \S\ref{sec:micro} requires.
  Simple agendas run above the diagonal (distrust forms first), complex ones below
  it (disagreement first), crossing over at $\alpha\approx1.7$.}
  \label{fig:trajectories}
\end{figure}

\paragraph{Why the agenda controls this.} The reason is exact, and it is a
conservation law. Every update to a weight vector is along $C_r\hat{x}$, and
$C_r\hat{x}$ always lies in the span of the agenda: it does initially, since
$C=I$, and the covariance update adds only multiples of
$(C\hat{x})(C\hat{x})^{\!\top}$, so it does at every step thereafter
(Appendix~\ref{app:conserved}). \emph{The component of each agent's weight vector
orthogonal to the span of the agenda is therefore conserved by the dynamics} --- in
the run of Figure~\ref{fig:polarisation}, $\|w_{\perp}\|=4.956$ at initialisation
and $4.956$ after $780\,000$ interactions.

Split $w_I = u_I + n_I$ into the learned part inside that span and the frozen part
outside it. Learning aligns the $u_I$ while the $n_I$ stay random and mutually
almost orthogonal, so
%
\begin{equation}
  \rho_{IJ} \;\approx\; \pm\, f_I f_J,
  \qquad f_I = \frac{\|u_I\|}{\|w_I\|},
  \label{eq:dilution}
\end{equation}
%
and opinion overlap is diluted by whatever fraction of each weight vector learning
cannot reach. Two facts then fix the ceiling. Because $\alpha<1$ leaves a
$(1-\alpha)$ fraction of directions untouched, $n_I\neq0$; and because the dynamics
anneals --- $C\to0$, so $\|u_I\|$ stops growing --- $f$ settles below one instead of
creeping towards it. Measured across four orders of magnitude in $\alpha$:

\begin{table}[h]
\centering\small
\begin{tabular}{lccccc}
\toprule
$\alpha = P/K$ & 0.03 & 0.17 & 0.50 & 1 & 3.3 \\
\midrule
$f^{2}$, predicted by \eqref{eq:dilution} & 0.29 & 0.62 & 0.84 & 1.00 & 1.00 \\
$\langle|\rho|\rangle$, measured            & 0.30 & 0.62 & 0.82 & 0.93 & 0.99 \\
\bottomrule
\end{tabular}
\end{table}

For $\alpha<1$ the prediction holds to $0.02$. At $\alpha\ge1$ the agenda spans the
whole space, the frozen component vanishes exactly, and $\langle|\rho|\rangle$ runs
to $0.99$; the residual gap at $\alpha=1$ is incomplete alignment rather than
geometry.

This is the ceiling behind Figure~\ref{fig:trajectories}, and it caps both
opinion-side quantities together: at $\alpha=0.17$ the society reaches
$\Rwmu=0.56$ and $\BI=0.19$, at $\alpha=3.3$ it reaches $0.93$ and $0.95$. The
difference is not that the narrow-agenda society has had less time --- no amount of
running will align directions that no issue ever probes. Trust has no such ceiling:
$\mu_{e|r}$ is one scalar per ordered pair, not a projection of a high-dimensional
vector, and $\eta=1-2\Phi(\hmu)$ saturates at $\pm1$ once $|\hmu|$ exceeds a few
units, at any $\alpha$. That asymmetry is why the trajectories climb the $\BA$ axis
first when the agenda is narrow.

\paragraph{This section is one phase of \S\ref{sec:results}.} With $d=0$ no class
information enters the dynamics anywhere, so the two class correlations vanish
identically --- measured as $\Rmuc=-0.004\pm0.030$ and $\Rcw=-0.003\pm0.017$ over
the whole range of $\fd$ --- while $\Rwmu$ alone is large. A society that has
polarised completely along a split unrelated to class is exactly region~(II) of the
phase diagram, the neutral phase. Everything above is region~(II) examined from the
inside, before the other three regions exist.

So the question of whether groups come to dislike each other because they disagree,
or come to disagree because they dislike each other, has a parameter attached to it
rather than an answer: narrow agendas --- few issues, functioning as symbols --- put
affect first, broad ones put opinion first, and the two are regimes of one system
rather than rival accounts. The same ratio explains, in \S\ref{sec:results}, why one
of the two discriminatory phases exists only for narrow agendas.
\section{Discrimination as an error of representation}
\label{sec:discrimination}

We now let some agents make a specific, minimal mistake: when they receive a message, they extend the
representation of the issue with a feature that carries no information about the
issue at all, and depends only on \emph{who sent it}. If that feature is
correlated with the emitter's class label, the agent is a
\emph{discriminating} agent, and the resulting shift of its opinion field is the
\emph{discrimination field}:
%
\begin{equation}
  \hw^{D} \;=\; \hw + D_{\kappa_r,\kappa_e},
  \label{eq:disc-field}
\end{equation}
%
where $D$ is a $2\times2$ matrix indexed by the classes of receiver and emitter.
A fraction $\fd$ of the society discriminates, drawn independently of class; the
rest use \eqref{eq:fields} unmodified.

This is a familiar failure viewed from an unusual angle: irrelevant features degrade
a classifier, and if one correlates with a protected attribute the errors become
structured along it. Here the classifier is one agent among many, and the question is
what those errors do to the \emph{population}.

\paragraph{What the field does.} By \eqref{eq:modulation}, the sign of perceived
agreement decides whether trust grows or decays. A receiver that adds $+d$ to
$\hw$ reads its counterpart as more agreeable than it is and becomes more
trusting: it is \emph{tolerant}. A receiver that adds $-d$ manufactures
disagreement that was not there and becomes more distrustful: it is
\emph{intolerant}. Throughout this paper $d>0$ therefore denotes
\textbf{tolerance towards in-group and intolerance towards out-group} emitters,
and $d<0$ the reverse --- discrimination in favour of the other class.

Geometrically, the field displaces the separatrix of the learning flow to
$\hmu = \hw + D$ (Figure~\ref{fig:flows}, Appendix~\ref{app:flows}). The undisplaced flow has two
symmetric consonant basins, ``trust and agree'' and ``distrust and disagree'',
and a society whose agents fall into them polarises into two camps whose
membership has nothing to do with class. Displacing the separatrix breaks that
symmetry: a receiver meeting an out-group emitter has an enlarged
``distrust and disagree'' basin, and one meeting an in-group emitter an enlarged
``trust and agree'' basin. Every discriminating agent is thus biased towards the
same alignment of distrust with class, and once enough of them are, the whole
society can lock into it.



\paragraph{Six ways to discriminate.} The matrix $D$ can be filled in six
inequivalent ways, listed in Appendix~\ref{app:cases}: one class may favour its
own without hostility to the other, be hostile without favouritism, or both; and
either one class or both may discriminate. All the phase diagrams below use the
fully symmetric case
%
\begin{equation}
  D \;=\; d\begin{pmatrix} +1 & -1\\ -1 & +1\end{pmatrix},
  \label{eq:case6}
\end{equation}
%
in which both classes favour their own and are hostile to the other. It is the
case in which the two classes play equivalent roles, so any asymmetry that
appears in the results is spontaneous rather than imposed.

\paragraph{A note on the sign convention.} An earlier version of this model
\citep{caticha2026discrimination} tabulates $D$ with the opposite sign while
placing the discriminatory phases at $d>0$ in its text and figures; the two are
inconsistent, and the inconsistency exchanges the halves of the phase diagram.
Appendix~\ref{app:sign} works this through and shows the maps under both
readings. We use \eqref{eq:disc-field} with \eqref{eq:case6} as stated, which is
the reading that agrees with the algorithm.

\subsection{Two more order parameters}
\label{sec:class-order-params}

Class now matters, so two statistics that refer to it can be added to those of
\S\ref{sec:order-params}. Writing $G_{IJ}=\kappa_I\kappa_J$ for the class
indicator ($+1$ for same-class pairs, $-1$ otherwise),
%
\begin{equation}
  \Rmuc = \bigl\langle G_{IJ}(\eta_{I|J}+\eta_{J|I}) \bigr\rangle,
  \qquad
  \Rcw  = \bigl\langle G_{IJ}\,\rho_{IJ} \bigr\rangle,
  \label{eq:correlations}
\end{equation}
%
again averaged over unordered pairs and normalised to $[-1,1]$. $\Rmuc$ is the order
parameter of discrimination: it asks whether trust follows class, distinguishing a
society that has merely polarised from one that has polarised \emph{along the
label}. $\Rcw$ asks whether opinion follows class, and separates the two
discriminatory phases from each other --- a society can have $\Rmuc$ near one with
$\Rcw$ near zero, and that combination is a distinct state rather than a degenerate
case. Appendix~\ref{app:conventions} records why $G$ is signed rather than
$\{0,1\}$ and how the three correlations are normalised.
\section{The phase diagram}
\label{sec:results}

We simulate societies of $N=40$ agents over a $200\times200$ grid in $(d,\fd)$, one
society per grid point, and measure the five order parameters of
\S\ref{sec:order-params} after a fixed number of interactions, with both classes
discriminating (case~\eqref{eq:case6}). Parameters and their provenance are in
Appendix~\ref{app:parameters}; each pixel is a single realisation, so the speckle
is society-to-society fluctuation, not a rendering artefact.

\begin{figure}[t]
  \centering
  \includegraphics[width=1.0\linewidth]{figures/order_parameter_maps}
  \caption{\textbf{All five order parameters over the $(d,\fd)$ plane}, for a
  simple agenda ($\alpha=0.17$, top) and a complex one ($\alpha=3.3$, bottom).
  $\Rmuc$ is the discrimination order parameter, flipping sign at $d\approx0$ ever
  more sharply as $\fd$ grows; $\Rcw$ is non-zero only for $d>0$ and, for the
  simple agenda, \emph{falls back} at the largest $d$, which is what separates the
  two discriminatory phases; $\BI$ separates the agenda sizes, featureless above
  and saturated below; and $\BA$ is the only panel in which the frustrated region
  at lower left announces itself, going strongly negative where everything else
  merely goes to zero.}
  \label{fig:correlations}
\end{figure}

\paragraph{Discrimination needs a quorum, and then arrives abruptly.} At small
$\fd$ the society stays neutral no matter how strong the bias: the few
discriminating agents are outvoted by the class-blind majority they learn from.
The trust--class correlation at large $d$ first exceeds $0.5$ around
$\fd\approx0.4$ ($0.43$ for the simple agenda, $0.35$ for the complex one), and past
that the transition in $d$ sharpens quickly: at $\fd=1$ the band in which
$|\Rmuc|<0.25$ is only $0.54$ wide, and $\Rmuc$ runs from $-0.99$ at $d=-1$ to
$+0.99$ at $d=+1$. This is the first practical consequence of the model: the state
of the population is not a smooth function of how biased its members are, and a
minority of biased agents leaves no trace at all until it reaches a threshold
fraction.

\paragraph{Two discriminatory phases, and only the simple agenda has both.} Both
discriminatory regions have distrust aligned to class; they differ in whether
opinion is aligned too. Along $\fd=0.9$ in the top row of
Figure~\ref{fig:correlations}, $\Rcw$ rises to $0.44$ at $d\approx0.6$ and then
\emph{falls} to $0.30$ by $d=1$, while $\Rmuc$ keeps climbing, from $0.78$ to
$0.87$ --- a gap of $0.55$ where the two started level. Opinion balance follows
$\Rcw$ down, $\BI$ dropping from $0.19$ to $0.06$ across the same row. At moderate bias the society splits into two camps
that both distrust and disagree with each other (region~III); at strong bias the
disagreement becomes unnecessary and class membership alone carries the hostility
(region~IV), which is what the magenta corner of Figure~\ref{fig:phase} shows.

The complex agenda has no region~IV. In the bottom row, $\Rmuc$ and $\Rcw$ track
each other across the whole $d>0$ half --- $0.84$ against $0.85$ at $d=0.15$, $0.91$
against $0.91$ at $d=0.75$, and still only $0.04$ apart at $d=1$ against the simple
agenda's $0.55$ --- while $\BI$ saturates at $0.98$. When there are many issues to disagree about, the
opinion split tracks the class split completely, so distrust and disagreement never
come apart. Class-only hostility requires the opposite: an agenda narrow enough that
most of every opinion lies in directions no issue probes, which by
\S\ref{sec:no-disc} caps $\Rcw$ however long the society runs. Region~(IV) is
therefore not a slow approach to region~(III) but a distinct state, available only
below $\alpha\approx1$.

\begin{figure}[t]
  \centering
  \includegraphics[width=0.55\linewidth]{figures/phase_diagram}
  \caption{\textbf{Four collective states.} The three correlations of
  Figure~\ref{fig:correlations} (simple agenda) composited as colour channels ---
  red $\Rmuc$, green $\Rcw$, blue $\Rwmu$ --- so that colour names the phase:
  (I) dark, frustrated reverse discrimination; (II) blue-violet, neutral
  polarisation unrelated to class; (III) pale, opinion tracking class; (IV) magenta,
  class alone.}
  \label{fig:phase}
\end{figure}

\paragraph{Favouring the out-group produces frustration, not inverted
discrimination.} For $d<0$ the correlations are the sign-mirror of the $d>0$ side,
which invites reading the left half of the diagram as ``reverse discrimination''.
The balance measures say otherwise (Figure~\ref{fig:correlations}). In the
discriminatory region $\BA=+0.97$: nearly every triple is balanced in trust, as
two coherent mutually hostile blocs require. At $d=-1$, $\fd=0.9$,
$\BA=-0.71$ for both agendas --- strongly \emph{anti}-balanced --- with
$\BI=0.00$, and $\Rwmu$ collapses to $0.03$ --- the only place it does. No
consistent assignment of trust exists: each agent is pulled towards the other class
while its neighbours are pulled towards theirs, so the society settles into a spin
glass rather than a mirror image. Pair correlations alone would report the region as
simply inverted; the triple-level measurement identifies it as disordered.
\section{Discussion}
\label{sec:discussion}

\paragraph{Optimality is context-dependent, and the context is group size.} The rule
these agents use is optimal for the problem it was derived for: one student, one
teacher, a noisy channel whose reliability must be inferred online. Nothing about it
is careless. Yet the same rule, run in a large population with one class-correlated
perturbation, converts an informationally empty label into a stable architecture of
group distrust. The mechanism is specific: it acts through $\Fmu$, and a rule with no
trust sector has none to bias --- under a plain Hebbian update the discrimination
field does nothing at all, since the update does not depend on $\hw$. We argue this
rather than measure it; the ablation is cheap and worth running. What exposes the
population is precisely the machinery that makes each agent good at inference, that
it draws conclusions about \emph{sources} as well as about the world. Giving
interacting agents a well-founded model of each other's reliability is not a neutral
improvement.

\paragraph{Individually small, collectively discontinuous.} A single discriminating
agent is biased by $O(d)$; the population's state is not, moving between phases over
a narrow interval in $d$ once $\fd$ passes its threshold. A protocol that inspects
agents one at a time --- the standard protocol --- can therefore certify every agent
as approximately unbiased while the population sits in a discriminatory phase.
Auditing a multi-agent system needs population-level order parameters, and the ones
used here cost two elicited matrices.

\paragraph{Two results that reframe standard questions.} The reversal at
$\alpha\approx1.7$ makes ``do groups disagree because they dislike each other, or the
reverse'' a question about a parameter rather than a choice between theories. And
making $d$ negative, so agents favour the out-group, gives not a tidy inverted
society but frustration:
$\BA<0$ with $\BI\approx0$, order parameters near zero not because nothing is
happening but because nothing can settle. Pair correlations alone would misreport
that state as neutral; only the balance measures identify it as a spin glass.

\looseness=-1
\paragraph{Limitations.} The agents are single-layer perceptrons over a fixed shared
embedding, all pairs may interact, and the society is closed and of fixed size. The
dynamics anneals rather than reaching a stationary state, so our phases are phases at
a finite interaction count, and the class labels are exactly inert --- the clean
version of a case that is never clean. Most importantly the result rests on one
substrate, agents whose learning rule we can write down; whether the same phases
appear where the agents are not analysable is the obvious next question, which the
order parameters of \S\ref{sec:order-params} were built to make answerable but which
we do not answer. Human societies are a further step again.

\subsubsection*{Reproducibility statement}
Every figure in this paper is produced by a script in the accompanying code,
which contains the model, the order parameters, and the sweep driver, together
with a test suite that checks the modulation functions against finite differences
of $\log Z$, the invariants of the dynamics (the covariance stays symmetric
positive-definite, the trust variance stays positive), and the order
parameters against their known limits. One command regenerates all figures at a
choice of three resolutions; sweeps are cached, so a figure can be restyled
without re-simulating.

The model has few parameters and we state all of them, including the ones the
formulation leaves open. \S\ref{sec:results} gives the values used;
Appendix~\ref{app:parameters} gives the full table, separating quantities fixed
by the model from those we chose, and records how each choice was made and what
happens under alternatives. Two conventions where a literal reading of the model
as previously written is ambiguous or inconsistent --- the sign of the
discrimination field, and the normalisation of the class indicator --- are
resolved explicitly in Appendices~\ref{app:sign} and \ref{app:conventions}, with
both readings implemented behind a flag so either can be reproduced.

\subsubsection*{Ethics statement}
This paper studies how discriminatory structure can arise among interacting
learning agents. Two points bear stating.

The model is not a claim about people. Its agents are single-layer perceptrons,
its groups are arbitrary inert labels, and its "tolerance" is a single scalar. We
use the vocabulary of in-group and out-group because it names the structure
compactly, not because the model licenses inferences about human prejudice, whose
causes are historical, institutional and material in ways nothing here
represents. Where we discuss the trust-versus-opinion polarisation
literature, we do so to point out that the model identifies a parameter that
selects between two regimes, not to adjudicate an empirical question about any
actual polity.

The practical direction of the work is defensive. The result implies that a
multi-agent system can occupy a discriminatory collective state while every agent
in it passes an individual audit, and that the transition into that state is
sharp in the strength of a per-agent bias. That argues for population-level
measurement of deployed multi-agent systems, and the order parameters used here
are inexpensive enough to serve as one.

\bibliography{references}
\bibliographystyle{iclr2027_conference}

\appendix
\section{The agenda's orthogonal complement is conserved}
\label{app:conserved}

Let $V_P=\mathrm{span}\{\hat{x}_1,\dots,\hat{x}_P\}$ be the subspace spanned by the
agenda. \emph{Claim:} for every agent, the component of $\hat{w}$ orthogonal to
$V_P$ is unchanged by \eqref{eq:update-ideo}.

By \eqref{eq:update-ideo} the weight update is a multiple of $C_r\hat{x}$ with
$\hat{x}\in V_P$, so it is enough that $C\hat{x}\in V_P$ at all times. At
initialisation $C=I$ and $C\hat{x}=\hat{x}\in V_P$. Suppose $C_t\hat{x}_i\in V_P$
for every issue. The covariance update is
$C_{t+1}=C_t+a\,(C_t\hat{x})(C_t\hat{x})^{\!\top}$ for the issue $\hat{x}$ just
seen, so for any issue $\hat{y}$,
%
\begin{equation}
  C_{t+1}\hat{y} \;=\; C_t\hat{y} \;+\; a\,(C_t\hat{x})\,\bigl(C_t\hat{x}\cdot\hat{y}\bigr),
\end{equation}
%
in which the first term lies in $V_P$ by hypothesis and the second is a multiple of
$C_t\hat{x}\in V_P$. Hence $C_t\hat{x}_i\in V_P$ for all $t$ by induction, every
increment to $\hat{w}$ lies in $V_P$, and the orthogonal component is conserved.
$\square$

Two consequences. For $\alpha=P/K<1$ each agent keeps a frozen random component of
its opinion that no interaction can touch, which dilutes every overlap by
\eqref{eq:dilution}; the test suite asserts the conservation numerically. For
$\alpha\ge1$ the agenda spans $\mathbb{R}^{K}$, the complement is empty, and the
ceiling disappears. Note that the conservation is a property of the \emph{agenda},
not of the discrimination field: it holds at every $d$ and $\fd$.

\section{All four modulation functions}
\label{app:modulation}

\begin{figure}[h]
  \centering
  \includegraphics[width=0.92\linewidth]{figures/modulation_contours_all}
  \caption{The two functions that move the means, $\Fw$ and $\Fmu$, above --- the
  pair shown in Figure~\ref{fig:modulation} --- and the two that anneal the
  uncertainties, $\FC$ and $\FV$, below. Each row shares a range. The mirror
  symmetry across $\hmu=\hw$ holds in both rows: $\Fw(x,y)=\Fmu(y,x)$ and
  $\FC(x,y)=\FV(y,x)$. $\FC$ and $\FV$ are negative almost everywhere, which is why
  both uncertainties shrink monotonically and the dynamics anneals
  (\S\ref{sec:micro}).}
  \label{fig:modulation-all}
\end{figure}

\section{Where blame for a surprise falls}
\label{app:slices}

\begin{figure}[h]
  \centering
  \includegraphics[width=0.95\linewidth]{figures/modulation_slices}
  \caption{Cuts through the modulation functions at fixed opinion field, $\hw=\pm6$.
  Below the crossover at $\hmu = \hw + D$ the trust sector moves (solid,
  $\Fmu$) and the receiver revises its trust; above it the opinion sector moves
  (dashed, $\Fw$) and the receiver unlearns what the emitter said. The
  discrimination field $D$ displaces the crossover --- dotted verticals, one per
  value of $D$ --- which is the entire mechanism studied in this paper, visible in a
  single coordinate.}
  \label{fig:modulation-slices}
\end{figure}

\section{The learning flow under a discrimination field}
\label{app:flows}

\begin{figure}[h]
  \centering
  \includegraphics[width=\linewidth]{figures/learning_flows}
  \caption{\textbf{The discrimination field bends the flow of learning.}
  Arrows are the change $(\Fw, \Fmu)$ that one interaction induces in the
  receiver's fields; the background is $|(\Fw,\Fmu)|$ and the thin contours are
  the evidence $Z$. Centre: without discrimination the flow escapes the two
  dissonant quadrants into two symmetric consonant basins, separated by
  $\hmu=\hw$ (green). Left and right: a discriminating receiver evaluates the
  flow at $\hw+D$, which moves the separatrix and enlarges one basin at the
  other's expense --- ``distrust and disagree'' for an out-group emitter
  ($D=-d$), ``trust and agree'' for an in-group one ($D=+d$).}
  \label{fig:flows}
\end{figure}

\section{Derivation of the learning algorithm}
\label{app:derivation}

The receiver's belief before an interaction is the product of Gaussians
\eqref{eq:belief}. It observes $\sigma_e$, the emitter's opinion on issue
$\hat{x}$, and treats it as a report of an unknown true label $\sigma_T$
transmitted through a channel that flips the label with probability
$\varepsilon(z)$. Marginalising the joint distribution of labels and of the issue
as the emitter may have parsed it, $y_t$,
%
\begin{equation}
  L(\sigma_e\mid \hat{x}, u) \;=\;
  \sum_{\sigma_T}\int P(\sigma_e\mid \sigma_T, y_t, \hat{x}, u)\,
  P(\sigma_T\mid y_t, \hat{x}, u)\,P(y_t\mid \hat{x}, u)\,\mathrm{d}y_t ,
\end{equation}
%
with $P(\sigma_e\mid\sigma_T,z)=(1-\varepsilon(z))\delta_{\sigma_e,\sigma_T} +
\varepsilon(z)\delta_{\sigma_e,-\sigma_T}$ and
$P(\sigma_T\mid y_t,w)=\Theta(\sigma_T\, y_t\cdot w)$. Two modelling choices keep
the algebra Gaussian. Taking $\varepsilon(z)=\Phi(z)$ with $z$ Gaussian of mean
$\mu_{e|r}$ and variance $V_{e|r}$ makes the flip probability a bounded function
of a Gaussian variable while preserving conjugacy of the trust sector --- a
Beta distribution for $\varepsilon$, the tempting alternative, does not. Taking
$P(y_t\mid \hat{x},u)$ Gaussian with variance $v_p=\|w\|^{-2}$ lets an
inexperienced agent, whose weights are small, be correspondingly unsure that it
has parsed the issue as the emitter did. Together these give the likelihood
%
\begin{equation}
  L(\sigma_e\mid z, \hat{x}, w) \;=\;
  \Phi(z)\,\Phi(-w\cdot\hat{x}\sigma_e)
  + \bigl(1-\Phi(z)\bigr)\,\Phi(w\cdot\hat{x}\sigma_e),
  \label{eq:likelihood}
\end{equation}
%
and the covariance stays block diagonal in the two sectors if it starts so.

Averaging the likelihood over the current belief gives the evidence
%
\begin{equation}
  Z(\sigma_e\mid \hat{x}, D_t) \;=\; \mathbb{E}_{Q_t}\bigl[L(\sigma_e\mid z, u, \hat{x})\bigr]
  \;=\;\sum_{\sigma_T}\mathbb{E}_z\bigl[P(\sigma_e\mid\sigma_T,\hat{x},z)\bigr]\,
  \mathbb{E}_w\bigl[\Phi(\sigma_T\, w\cdot\hat{x})\bigr],
\end{equation}
%
and the Gaussian integrals collapse onto the two scaled fields
\eqref{eq:fields}, giving \eqref{eq:evidence}: the evidence depends on the
receiver's state only through $(\hw,\hmu)$. This is the symmetry that makes the
two sectors mirror images.

The exact Bayesian posterior is not Gaussian. Replacing it by the Gaussian of
maximum entropy relative to the prior, subject to matching the first two moments
of the posterior, gives updates in terms of derivatives of the log evidence,
%
\begin{equation}
  \hat{w}_{t+1} = \hat{w}_t + C_t\frac{\partial \ln Z}{\partial \hat{w}},
  \qquad
  C_{t+1} = C_t + C_t \frac{\partial^{2}\ln Z}{\partial\hat{w}\,\partial\hat{w}^{\!\top}} C_t,
\end{equation}
%
and identically in the trust sector with $\mu_{e|r}, V_{e|r}$ in place of
$\hat{w}, C$. Expressing the derivatives in the scaled fields yields
\eqref{eq:update-ideo}--\eqref{eq:update-aff} with the modulation functions
\eqref{eq:modulation}. The first is a Hebbian term $\sigma_e\hat{x}$ with a
tensorial, self-annealing rate; the rest is the online estimation of the channel.

\section{Genealogy of the learning rule}
\label{app:genealogy}

The update \eqref{eq:update-ideo}--\eqref{eq:update-aff} descends from the
variational optimisation of modulated Hebbian on-line learning
\citep{kinouchi1992optimal,kinzel1990improving}, given a Bayesian reading by
\citet{opper1996online,opper1998bayesian} and applied across architectures by
\citet{solla1999optimal}; the entropic-dynamics formulation we use is
\citet{caticha2020entropic}. Learning through a noisy channel whose noise level
must itself be inferred was analysed by \citet{kinouchi1993lower} and
\citet{biehl1995learning}; \citet{alves2016distrust} identified that inferred
level with distrust and made it dynamical, with applications to voting
\citep{caticha2015whom} and to panels of appellate judges
\citep{caticha2019trust}, and a mean-field treatment by \citet{simoes2018mean}.
The blame-attribution behaviour resembles the least-action rule of
\citet{mitchison1989bounds}, which suffers retarded learning
\citep{kabashima1994perfect}, whereas the optimised algorithm attains
generalisation error decaying as $\alpha^{-1}$ even under output noise
\citep{simonetti1996online}.

\section{The six discrimination cases}
\label{app:cases}

The matrix $D$ in \eqref{eq:disc-field} admits six inequivalent fillings, where
$+d$ is the tolerant entry and $-d$ the intolerant one (\S\ref{sec:discrimination}
and Appendix~\ref{app:sign}). Rows index the receiver's class, columns the
emitter's.

\newcommand{\dmat}[4]{\left(\begin{smallmatrix}#1&#2\\#3&#4\end{smallmatrix}\right)}

\begin{table}[h]
\centering\small
\begin{tabular}{cccc}
\toprule
case & $D/d$ & who discriminates & how \\
\midrule
1 & $\dmat{+1}{0}{0}{0}$   & $A$ only & in-group favouritism \\
2 & $\dmat{0}{-1}{0}{0}$   & $A$ only & out-group hostility \\
3 & $\dmat{+1}{-1}{0}{0}$  & $A$ only & both \\
4 & $\dmat{+1}{0}{0}{+1}$  & both     & in-group favouritism \\
5 & $\dmat{0}{-1}{-1}{0}$  & both     & out-group hostility \\
6 & $\dmat{+1}{-1}{-1}{+1}$ & both    & both \\
\bottomrule
\end{tabular}
\end{table}

All phase diagrams in the main text use case 6, in which the two classes play
symmetric roles.

\section{The sign of the discrimination field}
\label{app:sign}

The sign of $D$ is fixed by the algorithm and not free. By \eqref{eq:modulation},
$\Fmu=(1-2\Phi(\hw))g(\hmu)/Z$ is negative when $\hw>0$, so perceived agreement
drives distrust $\mu$ \emph{down}. A receiver that adds $+d$ to $\hw$ therefore
becomes more trusting (tolerant) and one that adds $-d$ becomes more distrustful
(intolerant). Case 6 with $d>0$ consequently gives in-group trust and out-group
distrust, i.e. the discriminatory phase, which is the convention used throughout.

The companion manuscript \citep{caticha2026discrimination} tabulates the in-group
entry as $-d$ while its text and figures place the discriminatory phases at
$d>0$; those two statements are inconsistent, and taken literally the tabulated
signs put the discriminatory phases at $d<0$, mirroring every map in $d$. The same
inconsistency appears in its description of the learning flow, where the panel
labelled $D_{e|r}=-d$ is called the more tolerant one although $-d$ is the
intolerant entry, and where the ``distrust and disagree'' basin is said to grow
for $d>0$ although adding $+d$ shrinks it. A single global sign reconciles
everything; we adopt it in the form that leaves \eqref{eq:disc-field} untouched.
The accompanying code implements both readings behind a flag, and
Figure~\ref{fig:sign} runs the same sweep under each: the two rows are mirror
images in $d$, and only the first matches the maps as published.

\begin{figure}[h]
  \centering
  \includegraphics[width=0.92\linewidth]{figures/sign_convention_comparison}
  \caption{\textbf{The two readings of the discrimination field.} Trust--class
  correlation and trust balance over the $(d,\fd)$ plane under the convention
  used in this paper (top) and under the companion manuscript's table read
  literally with its Eq.~25 (bottom). The physics is identical; which half of the
  $d$ axis is discriminatory is not.}
  \label{fig:sign}
\end{figure}

\section{Conventions in the order parameters}
\label{app:conventions}

Two choices in \eqref{eq:correlations} deserve statement.

\textbf{The class indicator.} We use $G_{IJ}=\kappa_I\kappa_J\in\{-1,+1\}$. A
$\{0,1\}$ indicator, which counts only same-class pairs, cannot represent
anti-correlation: under perfect reverse discrimination it returns
$-(N/2-1)/(N-1)$ rather than $-1$, and its range depends on $N$. Since the
reverse-discrimination half of the phase diagram is a substantive part of the
result, the signed indicator is the appropriate one. Both are implemented.

\textbf{Normalisation.} $\Rwmu$ and $\Rmuc$ sum $\eta_{I|J}+\eta_{J|I}$ over
unordered pairs, so dividing by $N(N-1)$ places them in $[-1,1]$. $\Rcw$ has a
single term per pair and the same divisor would cap it at $1/2$; we use
$2/(N(N-1))$ so that all three share the range $[-1,1]$ and the three colour
channels of Figure~\ref{fig:phase} are commensurate.

\section{Closed form for the balance aggregates}
\label{app:closed-form}

Enumerating triples costs $O(N^{3})$ per society, which is prohibitive across a
phase diagram. For any matrix $M$ with unit diagonal --- either $\rho$ or $\eta$ ---
%
\begin{equation}
  \sum_{\substack{(I,J,K)\\ \text{distinct, ordered}}} M_{IJ}M_{JK}M_{KI}
  \;=\; \operatorname{tr}(M^{3}) - 3\Bigl(\textstyle\sum_{I,J}M_{IJ}M_{JI} - N\Bigr) - N ,
\end{equation}
%
and each unordered triple appears six times, three as the forward cycle and three
as the reverse, so
$\langle b\rangle = \bigl[\text{right-hand side}\bigr]/\bigl(6\binom{N}{3}\bigr)$.
For symmetric $M$ this is $\BI$; for the asymmetric trust matrix the two cycle
orientations are exactly the two terms of $b^{A}$, so the same expression gives
$\BA$. Both are batched matrix products. The identity is verified against
enumeration in the test suite.

\section{Parameters}
\label{app:parameters}

The model needs six numbers. One is recovered from the companion manuscript, one
is calibrated against its published figures, two are prior conventions, and two
are free choices; all six are listed with their provenance rather than left
implicit, since the manuscript states none of them. The accompanying code exposes
each of them and its calibration script reproduces the comparisons below.

\begin{table}[h]
\centering\small
\begin{tabular}{llll}
\toprule
symbol & meaning & value & provenance \\
\midrule
$K$ & dimension of the issue embedding & $30$ & recovered (see below) \\
$N$ & agents in the society & $40$ & chosen; converged (see below) \\
$P$ & issues on the agenda & $5$ / $100$ & chosen: $\alpha=P/K$ either side of $1$ \\
$\Delta t$ & interactions per ordered pair & $500$ & calibrated (see below) \\
$C_0, V_0$ & initial uncertainties & $I$, $1$ & uninformative prior \\
$\mu_0$ & initial distrust & $\mathcal{U}(-1,1)$ & gives $\BI\approx\BA\approx0$ at $t=0$ \\
\bottomrule
\end{tabular}
\end{table}

\paragraph{$K$ is recovered, not chosen.} The companion manuscript reports
balance trajectories at
$\alpha \in \{0.03, 0.17, 0.23, 0.33, 0.50, 0.67, 1.67, 3.33, 333.33\}$. These are
exactly $P/30$ for $P\in\{1,5,7,10,15,20,50,100,10^{4}\}$, which fixes $K=30$;
the language-agent protocol's thirty issues agrees.

\paragraph{$\Delta t$ is calibrated against the published trajectories.} The
dynamics anneals rather than reaching a stationary state, so the measurement time
matters and the companion manuscript does not state it (its text carries a
literal placeholder in place of the value). We scan $\Delta t$ and compare the
endpoints of all nine balance trajectories against the published ones:

\begin{table}[h]
\centering\small
\begin{tabular}{lccccc}
\toprule
$\Delta t$ (interactions per ordered pair) & 60 & 125 & 250 & \textbf{500} & 1000 \\
rms distance to published endpoints & 0.63 & 0.31 & 0.15 & \textbf{0.13} & 0.15 \\
\bottomrule
\end{tabular}
\end{table}

The minimum is flat-bottomed and sits at $\Delta t=500$, which we use throughout.
Seven of the nine endpoints then agree to within about $0.1$, and $\BI$ matches
the three largest agendas almost exactly, but no single $\Delta t$ reproduces all
nine: two middle curves ($\alpha=0.23$, $0.33$) sit above ours, and the two
largest agendas have a lower published $\BA$ than we obtain at any $\Delta t$ that
fits the rest. Either the society size differs from ours or those middle curves,
which overlap within a few pixels in the printed figure, cannot be digitised
reliably; we cannot distinguish the two. Dropping the ambiguous readings does not
move the selected value. All the $\Delta t$-independent checks pass: the
above/below-diagonal ordering with its crossover at $\alpha\approx1.7$, the sharp
sign flip of $\Rmuc$ at $d\approx0$, and the $\Rcw$ wedge confined to large $d$
and large $\fd$.

\paragraph{$N=40$ is large enough.} Repeating five representative points of the
plane at $N\in\{20,30,40,60\}$ with twelve realisations each, the order parameters
at $N=60$ differ from those at $N=40$ by at most $0.02$ ($\Rmuc$, $\Rcw$), $0.03$
($\BI$) and $0.05$ ($\BA$), while $N=20$ differs by up to $0.24$. The maps are
therefore converged in $N$ at $N=40$, and the cost of going further is steep: total
work grows as $N^{3}$.

\paragraph{Numerical safeguards.} Three floors are applied, none of which the
model prescribes: on the evidence $Z$, whose exact modulation functions diverge
in the deep dissonance corners; on the trust variance $V$, which $\FV<0$
drives monotonically towards zero; and on the covariance update, clipped to the
boundary at which $C+a(C\hat{x})(C\hat{x})^{\!\top}$ ceases to be positive
semi-definite. The last is the only one that can alter a trajectory, and it never
triggered in the runs reported here.

\end{document}
